%% kapa: exam=1 round=28 year=2017 period=1 form=A subject=civil_law subject_ko=민법 seq=29 number=29
%% source: Q-Net 감정평가사 기출문제 (cst00309) artlSeq=5199603 file=extracted/1008/5199603/2017년 제28회 감정평가사 1차시험 문제지/1교시 A형.pdf
%% conversion: kordoc 4.24.51 -> tools/gen1.py (auto); review=proofread-sonnet; flags=none
\begin{question}{29}
부동산 취득시효에 관한 설명으로 옳지 않은 것은? (다툼이 있으면 판례에 따름)
\begin{choices}
  \item 등기부취득시효의 요건으로서 무과실은 이를 주장하는 자가 증명하여야 한다.
  \item 점유취득시효에 있어서 점유자가 무효인 임대차계약에 따라 점유를 취득한 사실이 증명된 경우, 그 점유자의 소유의 의사는 추정되지 않는다.
  \item 시효취득자가 시효취득 당시 원인무효인 등기의 등기부상 소유명의자에게 시효이익을 포기한 경우에도 시효이익 포기의 효력이 발생한다.
  \item 점유취득시효의 완성 후 등기 전에 토지소유자가 파산선고를 받은 때에는 점유자는 파산관재인을 상대로 취득시효를 이유로 소유권이전등기를 청구할 수 없다.
  \item 토지에 대한 취득시효 완성으로 인한 소유권이전등기청구권은 그 토지에 대한 점유가 계속되는 한 시효로 소멸하지 않는다.
\end{choices}
\end{question}
\answer{3}
